\chapter{Mantenimiento y monitorización de aeromotores}\label{ch:22}
\objetivos{Vincular degradación, parámetros e inspección; construir una secuencia de diagnóstico que distinga síntoma, hipótesis y evidencia.}
\section{Una anomalía tiene contexto}
El motor se evalúa con configuración, historial de uso, ambiente, eventos y datos de mantenimiento. Un valor aislado puede variar por altitud, temperatura, régimen, carga y estado del sensor. La comparación de tendencias necesita condiciones comparables o correcciones justificadas. Una corrección empírica tampoco elimina toda dependencia ambiental.

La historia de arranques, ciclos térmicos, permanencias y excursiones ayuda a interpretar fatiga y fluencia. El número de horas no describe completamente esa historia. Se consideran piezas y módulos concretos y se conserva el seguimiento de límites de vida.
\section{Inspección por zonas}
En fan y compresor se examinan, según datos, daños por objetos extraños, erosión, suciedad, holguras y contactos. En cámara y turbina interesan grietas, quemaduras, deformación, pérdida de material o recubrimiento y otras indicaciones definidas. Las posiciones y dimensiones deben corresponder al sistema de referencia del manual.

La boroscopia permite observar zonas internas, pero una imagen oblicua puede distorsionar tamaño. La cobertura y la calidad se registran. Una indicación de aspecto similar puede tener criterios diferentes según etapa, borde o cara del álabe. No se inventa una tolerancia genérica de longitud de grieta.
\section{Lubricación, cojinetes y partículas}
El aceite lubrica, extrae calor y transporta partículas. Presión y temperatura se interpretan con régimen y sistema. Los filtros, detectores y análisis de aceite aportan información sobre desgaste y contaminación, pero requieren tendencia, método y origen probable. Una concentración puntual de metal no identifica de forma única un cojinete concreto.

Los cambios de aceite, horas desde toma y condiciones de muestreo afectan comparaciones. Una contaminación introducida al tomar la muestra puede generar una conclusión errónea. La evaluación combina indicaciones, inspección y datos del fabricante.
\section{Vibraciones y órdenes}
La vibración puede analizarse en tiempo, frecuencia y órdenes respecto a una velocidad de giro. La relación básica $f_n=nN/60$ conecta una orden $n$ con la velocidad $N$ en revoluciones por minuto. Un pico síncrono puede orientar una investigación de desequilibrio, pero no prueba por sí solo su causa. Alineación, holgura, rozamiento, excitación aerodinámica y montaje de sensor también intervienen.

La amplitud depende de dirección, punto, variable medida y procesamiento. Desplazamiento, velocidad y aceleración no se comparan como números equivalentes. Una resonancia puede amplificar respuesta en un intervalo de régimen; la evaluación no se reduce a escoger el pico mayor sin conocer condiciones.
\caso{M22. Orden de vibración}{Para un rotor ficticio a $12\,000\,\mathrm{rpm}$, la frecuencia de giro es $200\,\mathrm{Hz}$. Una componente a $400\,\mathrm{Hz}$ corresponde a orden 2 respecto a ese rotor. Si otro rotor gira a distinta velocidad, la misma frecuencia puede tener otra relación. El cálculo identifica una relación cinemática; no diagnostica un defecto ni define un límite de vibración.
\par\textbf{Desarrollo para comparar tu intento.} Una revolución por minuto equivale a $1/60$ ciclos por segundo. El orden es frecuencia observada dividida por frecuencia del rotor de referencia.

\begin{enumerate}

\item Convertir velocidad de giro a ciclos por segundo.

\[\fitmath{\mathrm{frot}=\frac{12000}{60}\approx 200\,\mathrm{Hz}}\]

\item Dividir la componente observada por la frecuencia del rotor.

\[\fitmath{\mathrm{order}=\frac{400}{\frac{12000}{60}}\approx 2\,}\]

\end{enumerate}

}
\section{Temperaturas y rendimiento}
La temperatura de escape, las velocidades de rotor, caudal de combustible y presiones se comparan bajo condiciones definidas. Un incremento de combustible para una prestación comparable puede indicar deterioro, pero también cambios de ambiente, configuración o medida. Los márgenes de temperatura y los criterios de retirada requieren la definición del fabricante.

La limpieza de compresor puede recuperar parte de una degradación por depósitos, pero no repara erosión, grietas o cambios permanentes de geometría. Su conveniencia y procedimiento dependen del motor y datos aplicables. Interpretar una recuperación de rendimiento como prueba de integridad sería incorrecto.
\ilustra{monitorizacion}{Cadena de diagnóstico: condiciones de medida, señal, hipótesis, corroboración e intervención documentada.}
\section{Diagnóstico y cierre}
Se define el síntoma, se valida la medición, se revisan eventos y se construyen hipótesis. Se priorizan pruebas por capacidad de discriminar, acceso y requisitos. Se contrasta evidencia y se decide dentro del procedimiento y alcance autorizado. La ausencia de confirmación de una hipótesis no prueba automáticamente la alternativa.

Tras intervenir, se comprueba el resultado con los datos y condiciones previstos. La documentación debe conservar el problema original, acciones y evidencias. Un cambio de componente «por prueba» puede no resolver la causa y requiere el control que proceda.
\section{Interpretación de tendencias con variables de confusión}
El margen EGT, las vibraciones y las partículas de aceite adquieren significado al compararse bajo condiciones operativas y métodos de normalización coherentes. Una serie que aumenta con horas de vuelo podría reflejar degradación, sesgo de sensor, cambio de motor, cambio de régimen o condiciones ambientales. Se evitan conclusiones causales a partir de correlación temporal aislada.

El orden vibratorio es una frecuencia normalizada por la velocidad de giro de su eje de referencia; un pico de primer orden no identifica por sí solo desequilibrio, y el análisis de órdenes exige conocer el eje y el método de adquisición. Los detectores de partículas y el análisis de aceite pueden apoyar hipótesis de desgaste, pero requieren confirmar fuente, contaminación y contexto. La boroscopia muestra únicamente la superficie accesible y observable. En una investigación se registran hipótesis rivales, señales discriminantes y evidencia antes y después de una intervención.
\repaso{¿Recuperar margen EGT después de una limpieza prueba que no había daño?}{No. La respuesta de rendimiento es compatible con ciertas causas, pero no excluye daño no observado y no reemplaza inspecciones exigidas.}


\section{Normalizar antes de comparar}
Temperatura, presión, velocidad y vibración cambian con régimen, ambiente y configuración. Comparar valores brutos puede confundir condición y degradación. Se usan correcciones y definiciones aplicables, conservando unidades y rotor de referencia. Se registran cambios de sensor, software, motor y componentes.

Un aumento de temperatura a prestación comparable puede sugerir pérdida de margen, sin identificar causa única. Se formulan hipótesis y buscan evidencias independientes. La tendencia detecta cambios; la causalidad requiere razonamiento físico.
\section{Frecuencia, orden y medida}
Orden es frecuencia dividida por giro del rotor de referencia, que se identifica en motores multieje. Seguir una orden o mantener frecuencia casi fija aporta indicios diferentes, sin diagnóstico automático. Amplitud, fase, posición, banda y condiciones también importan.

El muestreo y la ventana condicionan una señal digital; aliasing puede representar una frecuencia como otra. Se emplean equipos y procedimientos definidos, evitando reglas como «orden uno siempre significa desequilibrio». P33 verifica solo cinemática.
\section{Aceite y boroscopia como evidencias complementarias}
La muestra de aceite necesita identidad, método, punto y exposición. Concentración depende de aportación de material, reposición, filtrado y contaminación. Morfología de partículas orienta hipótesis; acciones y límites proceden de los datos pertinentes.

Boroscopia ofrece acceso y cobertura concretos. Iluminación, ángulo y referencia dimensional afectan medida. Lo no visible fuera del alcance del método no queda excluido. El libro no inventa límites de aceptación de álabe.
\section{Cerrar el ciclo de diagnóstico}
Se conserva señal, validación, hipótesis, comprobaciones, hallazgos y resultado de intervención. Si desaparece una señal, se comprueba el cambio de condiciones antes de atribuir causa. El rastro permite estudiar recurrencias y conectar daño material con decisiones documentadas.
\repaso{¿Dos señales que cambian juntas prueban una causa común?}{Pueden aumentar el interés de una hipótesis, pero también responder al régimen o a causas distintas. Se necesita evidencia discriminante.}

\resumen{Los parámetros necesitan condiciones y tendencia. Frecuencia y orden no son diagnóstico. La boroscopia tiene cobertura y límites de medida. Una recuperación de rendimiento no demuestra ausencia de daño. La decisión relaciona señal, mecanismo, datos de mantenimiento y evidencia de cierre.}
\fuentes{FAA--General, capítulos 5, 10 y 14 \cite{FAA-G}; CS-E 25 e información de mantenimiento \cite{CSE9}; documentos docentes de mantenimiento y mecanismos de materiales de la primera parte.}
